\section{Differential forms, pullbacks, and finite chains}\label{sec:geometry}

\subsection{Finite exterior algebra}

For a fixed ambient dimension $m$, a $q$-form is a finite coefficient array
\[
 \omega=\sum_{\substack{I=(i_1<\cdots<i_q)\\1\le i_j\le m}}
                   a_I\,\dd y_{i_1}\wedge\cdots\wedge\dd y_{i_q}.
\]
For $q=0$ the index list is empty and a form is a scalar coefficient.
Multiplication of basis words is defined as follows: concatenate their
index lists; if an index repeats, the product is zero; otherwise sort the
list and multiply by the sign $(-1)^{\text{number of inversions}}$.
This defines the wedge product by finite distributivity. Sorting,
duplicate detection, and inversion counting are all bounded operations.

For $C^1$ coefficients, the exterior derivative is defined by
\[
 \dd\omega=\sum_I\sum_{a=1}^m
          \partial_a a_I\,\dd y_a\wedge\dd y_{i_1}\wedge\cdots\wedge\dd y_{i_q}.
\]
Equivalently, its coefficient at $K=(k_1<\cdots<k_{q+1})$ is
\[
 (\dd\omega)_K
   =\sum_{\nu=1}^{q+1}(-1)^{\nu-1}
        \partial_{k_\nu}a_{K\setminus k_\nu}.
\]
These formulas specify all signs without a topological convention left
implicit.

Finite expansion and the scalar product rule prove the graded Leibniz
identity
\[
 \dd(\alpha\wedge\beta)
 =\dd\alpha\wedge\beta+(-1)^{\deg\alpha}\alpha\wedge\dd\beta.
\]
Indeed, differentiating a coefficient product produces its two scalar
terms. In the second term the new one-form must move past
$\deg\alpha$ basis one-forms, giving the displayed sign. Repeated-index
terms vanish in both expansions. This is a finite matching of coefficient
terms.

\subsection{The pullback construction}

Let $\sigma:Q^k\to D\subseteq\R^m$ be a certified parameterization.
Define
\[
 \dd\sigma_a=\sum_{j=1}^k\partial_j\sigma_a\,\dd x_j,\qquad
 \sigma^*\omega=\sum_I(a_I\circ\sigma)
                  \dd\sigma_{i_1}\wedge\cdots\wedge\dd\sigma_{i_q}.
\]
For an increasing parameter index list $J$ of length $q$, the coefficient
is
\begin{equation}\label{eq:pullback-minors}
 (\sigma^*\omega)_J
 =\sum_{|I|=q}(a_I\circ\sigma)
        \det\bigl(\partial_{j_b}\sigma_{i_a}\bigr)_{1\le a,b\le q}.
\end{equation}
The determinant here means its permutation expansion. Its construction
requires no inverse matrix or rank decision.

If $\omega$ has certified $C^1$ coefficients and $\sigma$ has a certified
$C^2$ jet, every factor in \eqref{eq:pullback-minors} has a $C^1$
presentation: use the composition lemma for $a_I\circ\sigma$, and the
supplied second jets for $\partial_j\sigma_a$. Apply the sum and product
lemmas finitely many times. This constructs the complete $C^1$ certificate
for the pullback.

For explicit size bounds, suppose a common integer $M\ge1$ bounds all
ambient coefficient values and first derivatives, and all first and
second derivatives of $\sigma$. Then
\[
 \binom mq q!M^{q+1}
\]
bounds every pullback coefficient. A safe bound on any of its first
derivatives is
\[
 \binom mq q!(m+q+1)M^{q+2}.
\]
To see the latter, a derivative falling on $a_I\circ\sigma$ gives at most
$m$ terms; a derivative falling on a determinant gives at most $q$ terms
per permutation. Each term is a product of the bounded factors. This
bound also covers $q=0$. The effective modulus construction is the finite
expression-tree recursion of \cref{sec:product,sec:chain-rule}.

\begin{lemma}[Naturality of the exterior derivative]\label{lem:naturality}
For these certified presentations,
\[
 \dd(\sigma^*\omega)=\sigma^*(\dd\omega)
\]
as coefficientwise equality of named forms.
\end{lemma}
\begin{proof}
First,
\[
 \dd(\dd\sigma_a)
 =\sum_{i<j}
   \bigl(\partial_i\partial_j\sigma_a
          -\partial_j\partial_i\sigma_a\bigr)\dd x_i\wedge\dd x_j=0
\]
by the symmetry certificates for the second jet. For one summand of
$\sigma^*\omega$, apply the graded product rule:
\begin{align*}
 &\dd\bigl((a_I\circ\sigma)
                  \dd\sigma_{i_1}\wedge\cdots\wedge\dd\sigma_{i_q}\bigr)\\
 &\quad=\dd(a_I\circ\sigma)
                  \wedge\dd\sigma_{i_1}\wedge\cdots\wedge\dd\sigma_{i_q}.
\end{align*}
All the other terms contain a factor $\dd(\dd\sigma_{i_j})$ and vanish.
By the chain rule already proved,
$\dd(a_I\circ\sigma)=\sum_a(\partial_a a_I\circ\sigma)\dd\sigma_a$.
Substitution gives exactly the pullback of
$\sum_a\partial_a a_I\,\dd y_a\wedge\dd y_I$. Sum over $I$.
Every equality is a finite coefficient expansion, using the explicit
jet closure lemmas, and no integration theorem has entered the proof.
\end{proof}

\subsection{Integrating a parameterized cell}

For a $k$-form $\eta$ with certified uniformly continuous coefficients,
define
\[
 \int_\sigma\eta:=\int_{Q^k}\sigma^*\eta.
\]
The right side is the named integral of its single top-degree coefficient.
For $\eta=\dd\omega$, the supplied first-derivative continuity of
$\omega$ provides this condition; a second derivative of the ambient
form is not required.
The boundary of $\sigma$ is the formal chain
\begin{equation}\label{eq:sigma-boundary}
 \partial\sigma
 =\sum_{i=1}^k(-1)^{i-1}
         \bigl(\sigma\circ\delta_i^1-\sigma\circ\delta_i^0\bigr).
\end{equation}
For a zero-dimensional cell, integration of a zero-form is its value
at the point, and its boundary is zero.

Restriction of a jet to a face merely deletes the frozen coordinate
directions; its moduli and bounds can be retained. Direct substitution
in the finite minor formula proves
\[
 (\sigma\circ\delta_i^\varepsilon)^*\omega
       =(\delta_i^\varepsilon)^*(\sigma^*\omega).
\]
Thus the boundary integral defined from \eqref{eq:sigma-boundary} is the
boundary integral of the pulled-back form, with exactly the signs used
in the cube theorem.

\begin{theorem}[Stokes for a certified curved cube]\label{thm:curved}
Let $\sigma$ be a certified $C^2$ $k$-cube in a rational ambient box, and
let $\omega$ be a certified $C^1$ $(k-1)$-form there. Then
\[
 \int_\sigma\dd\omega\ \EqName\ \int_{\partial\sigma}\omega .
\]
The pullback construction and the schedule
\eqref{eq:schedule-t}--\eqref{eq:schedule-values} give primitive recursive
names and explicit equality witnesses.
\end{theorem}
\begin{proof}
The effective calculus constructs the $C^1$ coefficient presentation of
$\alpha=\sigma^*\omega$. Write it in flux form \eqref{eq:flux-form}.
Its diagonal derivative data satisfy (D), and all coefficients and
diagonal derivatives satisfy (C) and (U). Apply \cref{thm:rates}.
By \cref{lem:naturality}, the volume side is
$\int_{Q^k}\sigma^*(\dd\omega)=\int_\sigma\dd\omega$.
The face-restriction identity identifies the boundary side with
$\int_{\partial\sigma}\omega$.
\end{proof}

\subsection{Finite chains and the error budget}

Define the integral over $C=\sum_{\lambda<L}c_\lambda\sigma_\lambda$
by the same finite integer linear combination of integrals, and define
$\partial C=\sum c_\lambda\partial\sigma_\lambda$. Put
$W=\sum_{\lambda<L}|c_\lambda|$. If $W=0$, both sides are zero.
Otherwise let $s=\ell(W)$ and define
\[
 V_C(r)=\sum_{\lambda<L}c_\lambda V_\lambda(r+s),\qquad
 B_C(r)=\sum_{\lambda<L}c_\lambda B_\lambda(r+s).
\]
Here the individual names are those in \cref{thm:curved}. The finite
triangle inequality multiplies every bound by at most $W$, and
$W2^{-s}\le1$. Hence the Cauchy and equality rates of
\cref{thm:rates} remain valid for $V_C,B_C$. This proves
\cref{thm:informal}. One can use separate meshes for separate cells,
or take the finite maximum of their moduli and use a common mesh.

\subsection{Why the boundary has no boundary}

Faces can be represented by the original cell identifier and a sorted
list of frozen original coordinates with their endpoint values.
This representation makes two successive face restrictions normalize
to the same data when they freeze the same two coordinates.

Fix original coordinates $i<j$ and endpoint values
$\varepsilon,\eta$. Freezing $i$ first makes the later index for $j$
equal to $j-1$. Its sign in the double boundary is
\[
 (-1)^{i-1+j-2}(2\varepsilon-1)(2\eta-1).
\]
Freezing $j$ first leaves the later index for $i$ equal to $i$.
The sign is
\[
 (-1)^{j-1+i-1}(2\eta-1)(2\varepsilon-1),
\]
which is its negative. Each codimension-two face occurs exactly in
these two ways. Their coefficients cancel, proving
\[
 \partial^2C=0
\]
by a finite sign-reversing pairing. No decision procedure for extensional
equality of arbitrary parameterization programs is required.

\subsection{Matching faces and orientation changes}

For different cells, the fact that two geometric faces coincide is not
decided by comparing arbitrary programs. An interface matching record
may instead supply an explicit equality of their parameterizations and
tangential first jets, with named equality witnesses. The continuity
and finite-expression bounds then prove equality of their pulled-back
coefficients. If their signed multiplicities are opposite, their
integrals cancel.

A useful supported reparameterization is a permutation and reflection
of cube coordinates. Let $T:Q^d\to Q^d$ be such a signed coordinate
permutation, with constant determinant $s\in\{-1,1\}$. The midpoint
grid $\{(a+\tfrac12)/N\}$ is permuted exactly by $T$. For a top form
$b(x)\,\dd x_1\wedge\cdots\wedge\dd x_d$, the rational midpoint sums
for its pullback, using the direct valid evaluator
$H_T(x,p)=sH(Tx,p)$, are therefore exactly $s$ times the original sums.
The general guarded composition evaluator is equivalent to this one,
although it need not return the same rational sample at the same
precision. Uniform coefficient equality identifies their integral names.
The midpoint and left-grid sums are within the same (U) bound,
because corresponding tags are at distance at most $1/N$.
Their names agree by refinement. Consequently
\[
 \int_{Q^d}T^*\eta=s\int_{Q^d}\eta.
\]
This proves the orientation rule needed for matching faces after a
coordinate symmetry. A matching record must include this sign.

Dyadic subdivision also preserves an integral: use direct rational affine
evaluation and the exact determinant weight on each subcube. The
resulting points and weights are exactly those in the refined whole-cube
sum. Finite grouping proves the identity, and the guarded evaluator
represents the same named integral. Thus many
presentation changes can be certified without a general
change-of-variables theorem. Invariance under an arbitrary smooth
diffeomorphism is a separate theorem and is not silently assumed here.

\section{Classical vector-calculus forms}\label{sec:vector}

\subsection{Kelvin--Stokes on a surface patch}

Let $\sigma(u,v)=(\sigma_1,\sigma_2,\sigma_3)$ be a certified $C^2$ patch,
and let $\mathbf F=(F_1,F_2,F_3)$ have a certified $C^1$ jet.
For $\omega=F_1\dd x+F_2\dd y+F_3\dd z$ set
\[
 P=\mathbf F(\sigma)\cdot\sigma_u,\qquad
 Q=\mathbf F(\sigma)\cdot\sigma_v.
\]
Then $\sigma^*\omega=P\,\dd u+Q\,\dd v$. The flux coefficients of
the cube theorem are $f_1=Q$ and $f_2=-P$.

The product and chain rules give
\begin{align*}
 \partial_uQ-\partial_vP
 &=\sum_{a,b}(\partial_aF_b)(\sigma)
       \bigl((\partial_u\sigma_a)(\partial_v\sigma_b)
             -(\partial_v\sigma_a)(\partial_u\sigma_b)\bigr).
\end{align*}
Terms containing $F_b(\sigma)\partial_u\partial_v\sigma_b$ cancel
by the second-jet symmetry. Expanding the three-component cross
product shows that the remaining finite polynomial is
\[
 (\curl\mathbf F)(\sigma)\cdot(\sigma_u\times\sigma_v).
\]
Therefore \cref{thm:curved} becomes
\begin{align}
 &\int_0^1P(u,0)\,\dd u+\int_0^1Q(1,v)\,\dd v
       -\int_0^1P(u,1)\,\dd u-\int_0^1Q(0,v)\,\dd v\notag\\
 &\hspace{2em}=
 \int_{Q^2}(\curl\mathbf F)(\sigma(u,v))
                       \cdot(\sigma_u\times\sigma_v)\,\dd u\,\dd v .
 \label{eq:kelvin-explicit}
\end{align}
All equalities in this section are equalities of the constructed names.
The boundary runs counterclockwise in the parameter square.

For an embedded oriented surface given by compatible patches, the right
side is the usual oriented surface integral, and the matched interior
edges cancel. The vector area element is used directly:
$\mathbf n\,\dd S=(\sigma_u\times\sigma_v)\dd u\,\dd v$.
There is no need to divide by $|\sigma_u\times\sigma_v|$ or compute a unit
normal. An immersion certificate is needed only for the additional
geometric interpretation as a regular embedded patch, not for the
parameterized identity itself.

\subsection{Green, divergence, and exact-form consequences}

For a planar patch, $\omega=P\,\dd x+Q\,\dd y$ gives
\[
 \int_{\partial C}P\,\dd x+Q\,\dd y
   =\int_C(\partial_xQ-\partial_yP)\,\dd x\wedge\dd y.
\]
For a three-dimensional chain use
\[
 \omega=F_1\,\dd y\wedge\dd z
       +F_2\,\dd z\wedge\dd x+F_3\,\dd x\wedge\dd y.
\]
Its exterior derivative is
$(\partial_xF_1+\partial_yF_2+\partial_zF_3)\dd x\wedge\dd y\wedge\dd z$,
giving the divergence theorem for the supplied presentation.
For $k=1$, the theorem is the fundamental theorem of calculus for
certified parameterized paths; only $C^1$ path data are needed in that
case.

If $\partial C=0$, then $\int_C\dd\eta=0$ whenever the indicated
certificates are available. If $C_1-C_2=\partial D$ is a \emph{supplied}
finite filling and $\dd\omega=0$ with a certificate, then
$\int_{C_1}\omega=\int_{C_2}\omega$.
These are constructive consequences of finite chain algebra and the
proved Stokes identity. They do not assert the existence of a filling
or solve a topological recognition problem.
